\documentclass[12pt]{article}
\usepackage{amsmath,amssymb,amsfonts}
\usepackage[margin=1in]{geometry}

\begin{document}

\section*{Problem 5.26}

\textbf{Problem:} Let $\hat{X}$ be the self-adjoint operator on the Hilbert space of square-integrable functions defined by $\hat{X}f(x) = xf(x)$.
\begin{itemize}
\item[(i)] Is $\hat{X}$ well defined on all elements of the space? If not, characterize the elements for which it is defined.
\item[(ii)] Prove that $E_\xi(\hat{X})$ defined by
\[
E_\xi(\hat{X})f(x) = \begin{cases} f(x) & \text{if } x \leq \xi, \\ 0 & \text{if } x > \xi, \end{cases}
\]
is the resolution of the identity belonging to $\hat{X}$.
\end{itemize}

\bigskip
\textbf{Solution:}

\subsection*{Part (i): Domain of $\hat{X}$}

The operator $\hat{X}$ maps $f(x) \mapsto xf(x)$. For $\hat{X}f$ to be in $L^2$, we need:
\[
\|\hat{X}f\|^2 = \int_{-\infty}^{\infty}|xf(x)|^2\,dx = \int_{-\infty}^{\infty}x^2|f(x)|^2\,dx < \infty.
\]

This is \emph{not} guaranteed for all $f \in L^2$. For example, $f(x) = (1+x^2)^{-1/2}$ is in $L^2$ but $xf(x) = x/(1+x^2)^{1/2}$ is not in $L^2$.

\[
\boxed{\text{The domain of } \hat{X} \text{ is } D(\hat{X}) = \left\{f \in L^2 : \int_{-\infty}^{\infty}x^2|f(x)|^2\,dx < \infty\right\}.}
\]

This is a dense subspace of $L^2$ but not the whole space. $\hat{X}$ is an unbounded operator.

\subsection*{Part (ii): $E_\xi$ is the spectral resolution of $\hat{X}$}

We need to verify that $E_\xi$ satisfies the properties of a spectral family and that the spectral theorem gives back $\hat{X}$:
\[
\hat{X} = \int_{-\infty}^{\infty}\lambda\,dE_\lambda.
\]

\textbf{Property 1: $E_\xi$ is a projection.}
\[
E_\xi^2 f(x) = E_\xi(E_\xi f)(x) = \begin{cases} E_\xi f(x) = f(x) & x \leq \xi, \\ 0 & x > \xi. \end{cases} = E_\xi f(x). \quad\checkmark
\]

\textbf{Property 2: $E_\xi$ is self-adjoint.}
\[
(E_\xi f, g) = \int_{-\infty}^{\xi}f(x)\overline{g(x)}\,dx = (f, E_\xi g). \quad\checkmark
\]

\textbf{Property 3: Monotonicity.} If $\xi_1 < \xi_2$, then $E_{\xi_1}f(x) = f(x)\chi_{(-\infty,\xi_1]}(x)$, and for any $f$:
\[
\|E_{\xi_1}f\|^2 = \int_{-\infty}^{\xi_1}|f|^2\,dx \leq \int_{-\infty}^{\xi_2}|f|^2\,dx = \|E_{\xi_2}f\|^2.
\]
So $E_{\xi_1} \leq E_{\xi_2}$. $\checkmark$

\textbf{Property 4: Limits.}
\[
\lim_{\xi\to-\infty}E_\xi = 0, \qquad \lim_{\xi\to+\infty}E_\xi = I. \quad\checkmark
\]

\textbf{Property 5: Recovery of $\hat{X}$.}

For the spectral integral, $dE_\lambda$ acts as multiplication by $\chi_{(\lambda, \lambda+d\lambda]}(x)$:
\[
\left(\int_{-\infty}^{\infty}\lambda\,dE_\lambda\right)f(x) = \int_{-\infty}^{\infty}\lambda\,\chi_{(\lambda,\lambda+d\lambda]}(x)\,f(x) = xf(x) = \hat{X}f(x).
\]

More rigorously:
\[
(f, \hat{X}g) = \int_{-\infty}^{\infty}x\overline{f(x)}g(x)\,dx = \int_{-\infty}^{\infty}\lambda\,d(f, E_\lambda g) = \int_{-\infty}^{\infty}\lambda\,d\left(\int_{-\infty}^{\lambda}\overline{f(x)}g(x)\,dx\right).
\]

Since $\frac{d}{d\lambda}\int_{-\infty}^{\lambda}\overline{f}g\,dx = \overline{f(\lambda)}g(\lambda)$, the Stieltjes integral gives:
\[
\int_{-\infty}^{\infty}\lambda\,\overline{f(\lambda)}g(\lambda)\,d\lambda = \int_{-\infty}^{\infty}\lambda\,\overline{f(\lambda)}g(\lambda)\,d\lambda = (f, \hat{X}g). \quad\checkmark
\]

\[
\boxed{E_\xi(\hat{X})f(x) = \chi_{(-\infty,\xi]}(x)\,f(x) \text{ is the spectral resolution of } \hat{X} = x.}
\]

\end{document}
